% from ega4-2-fr.tex
\oldpage[IV-2]{32}
\emph{(xv)} Si $f'$ est fini, il est affine, donc aussi $f$ d'après (xiii)~;
en outre, avec les notations de la démonstration de (xiii), $\mathcal A'$ est
un $\mathcal O_{Y'}$-Module de type fini, et $\mathcal A'$ est isomorphe à
$g^*(\mathcal A)$~; il résulte de (2.5.2) que $\mathcal A$ est un
$\mathcal O_Y$-Module de type fini, donc $f$ est un morphisme fini.

\emph{(xvi)} Dire que $f$ est quasi-fini signifie que $f$ est un morphisme de
type fini et que pour tout $y\in Y$, $f^{-1}(y)$ est fini (II, 6.2.2 et I,
6.4.4)~; la conclusion résulte donc de (v) et de (xv).

\emph{(xvii)} On voit comme dans (xv) que $f$ est affine. On peut se borner
au cas où $Y=\operatorname{Spec}(A)$, $Y'=\operatorname{Spec}(A')$, et alors
$X=\operatorname{Spec}(B)$, $X'=\operatorname{Spec}(B')$, où
$B'=B\otimes_A A'$~; $B$ est égale à la limite inductive de ses
sous-$A$-algèbres de type fini $B_\alpha$, donc on a
$B'=\varinjlim B'_\alpha$, où $B'_\alpha=B_\alpha\otimes_A A'$, et
$B'_\alpha$ est une $A'$-algèbre de type fini. Mais par hypothèse $B'$ est
entière sur $A'$, donc $B'_\alpha$ est un $A'$-module de type fini, et
$B_\alpha$ est donc un $A$-module de type fini (2.5.2). C.Q.F.D.

\begin{corollary}[2.7.2]
\label{IV.2.7.2-fr}
Les hypothèses et notations étant celles de (2.7.1), supposons $f$
quasi-compact~; soient $\mathcal L$ un $\mathcal O_X$-Module inversible,
$\mathcal L'=\mathcal L\otimes_{\mathcal O_X}\mathcal O_{X'}$, son image
réciproque. Pour que $\mathcal L$ soit ample (resp. très ample) pour $f$, il
faut et il suffit que $\mathcal L'$ soit ample (resp. très ample) pour $f'$.
\end{corollary}

La condition est nécessaire sans hypothèse sur $g:S'\to S$ (II, 4.4.10 et
4.6.13)~; pour voir qu'elle est suffisante, on peut comme dans (2.7.1) se
limiter au cas où $S=Y$, $S'=Y'$. L'hypothèse sur $\mathcal L'$ implique que
$f'$ est quasi-compact et séparé (II, 4.6.1), donc il en est de même de $f$
((2.6.2, (v)) et (2.7.1, (i))). Posons
$\mathcal E=f_*(\mathcal L)$, $\mathcal E'=f'_*(\mathcal L')$~; il résulte
de (2.3.1) que l'homomorphisme canonique
$u:g^*(\mathcal E)\to\mathcal E'$ est bijectif. Si $\mathcal L'$ est très
ample pour $f'$, l'homomorphisme canonique
$\sigma':f'^*(\mathcal E')\to\mathcal L'$ est surjectif, et le morphisme
\[
r'=r_{\mathcal L',\sigma'}:X'\longrightarrow\mathbf P(\mathcal E')
\]
est une immersion (II, 4.4.4, b)) nécessairement quasi-compacte (1.1.2,
(v)). Le fait que $u:g^*(\mathcal E)\to\mathcal E'$ soit bijectif entraîne
que si $h:\mathbf P(\mathcal E)\to Y$, $h':\mathbf P(\mathcal E')\to Y'$
sont les morphismes structuraux, $h'$ s'identifie à $h_{(Y')}$ (II, 4.1.3)~;
d'autre part, en désignant par $g'$ la projection $X'\to X$, $g'$ est
fidèlement plat (2.2.13), on a $f\circ g'=g\circ f'$, et l'on vérifie
aisément que l'homomorphisme
\[
g'^*(f^*(\mathcal E))\longrightarrow g'^*(\mathcal L)
\]
est identique à l'homomorphisme composé
\[
f'^*(g^*(f_*(\mathcal L)))\xrightarrow{f'^*(u)}
f'^*(f'_*(\mathcal L'))\xrightarrow{\sigma'}\mathcal L'
\]
(par exemple en se ramenant au cas où $Y$ et $Y'$ sont affines). Comme
$\sigma'$ est surjectif et $f'^*(u)$ bijectif, on voit que
$g'^*(\sigma)$ est surjectif, donc il en est de même de $\sigma$ (2.2.7).
On en conclut que le morphisme
$r=r_{\mathcal L,\sigma}:X\to\mathbf P(\mathcal E)$ est partout défini
(II, 3.7.4)~; en outre, si l'on pose $P=\mathbf P(\mathcal E)$,
$P'=\mathbf P(\mathcal E')$ et si $g''$ est la projection $P'\to P$, $r'$
s'identifie à $r_{(P')}$ (II, 4.2.10) et $g''$ est fidèlement plat et
quasi-compact (1.1.2 et 2.2.13). On conclut donc de (2.7.1, (xi)) que $r$ est
une immersion, et par suite $\mathcal L$ est très ample (II, 4.4.4, b)).

Supposons maintenant que $\mathcal L'$ soit ample pour $f'$~; pour démontrer
que $\mathcal L$ est ample pour $f$, on peut se borner au cas où $Y$ est
affine (II, 4.6.4), et en vertu de (2.2.12) et de (II, 4.6.13), on peut aussi
supposer que $Y'$ est affine. Alors $X$ et $X'$ sont des

